\chapter{Computable Function Operators}
\label{sec:p04-function-operators}

Computing with a function requires a name of the function, not
only the assertion that each of its values exists.
This chapter develops evaluation, integration, differentiation and
zero finding in the name/realizer framework of
\cref{sec:p04-represented-spaces}.
Its primary source is the corrected second version, July 1995,
of Weihrauch's \emph{A Simple Introduction to Computable Analysis},
Chapters 6--7, pages 42--54 \cite{Weihrauch1995IntroductionRS}.
Its source-native representations and operand types are retained
below. Local reconstructions supply proofs where that text
explicitly omits one.

\section{Names of continuous functions}
\label{sec:p04-fo-names}

\begin{hbconvention}{def:p04-fo-native-notation}{Source alphabet, rational words and typed operands}
In the July 1995 presentation, \(\Sigma\) is a sufficiently
large finite alphabet, \(\Sigma^*\) consists of finite words
and \(\Sigma^\omega\) of infinite symbol streams.
Its separators are \(\#\) and the distinct cent character
\(\text{\textcent}\).
The notation \(\nu_{\mathrm{bin}}\) reads the word \(0\)
or a binary word beginning with \(1\) as a natural number.
The native \(\nu_{\mathbb Q}:\Sigma^*\rightharpoonup\mathbb Q\)
reads such nonnegative integer words, their nonzero negatives,
and signed reduced fractions \(u/v\), with nonzero numerator,
denominator greater than one, and relatively prime numerator
and denominator. It is undefined on other finite words
(Appendix C, Definition C1).
Write \(\bar u=\nu_{\mathbb Q}(u)\) on its domain.
This finite-word notation and the host numerical rational
decoder in \cref{def:p04-rs-rational-code} have different
input types; the shared glyph does not identify them.

Write \(C[0,1]\) for real-valued continuous functions on
\([0,1]\), \(C(\mathbb R)\) for those on \(\mathbb R\), and
\(\mathrm{Pg}\) for functions linear between successive
rational vertices, with domain endpoints zero and one.
Define \(\lVert h\rVert_\infty
=\max_{0\leq t\leq1}|h(t)|\) and
\(d(f,g)=\lVert f-g\rVert_\infty\).
Throughout the native clauses below, \(\rho_C\) names a
scalar real, \(\delta_C\) names a function on \([0,1]\),
and \(\delta_{\mathbb R}\) names a function on \(\mathbb R\).
In contrast, host \(\delta_{\mathrm C}\) consumes a Baire
stream of numerical rational codes and names a scalar real.
Every operand's domain and codomain, not subscript typography,
determines the representation being used.
\end{hbconvention}

\begin{hbdefinition}{def:p04-fo-native-scalar}{Native scalar Cauchy names: \(\rho_C\)}
The source Definition 3.1 defines
\(\rho_C:\Sigma^\omega\rightharpoonup\mathbb R\) by
\[
 \rho_C(p)=x
 \quad\text{iff}\quad
 p=u_0\#u_1\#\cdots,\quad
 (\forall n)(\forall j>n)\ |\bar u_j-\bar u_n|<2^{-n},
 \quad x=\lim_n\bar u_n,
\]
with every \(u_n\in\dom(\nu_{\mathbb Q})\).
The pairwise clause implies
\(|\bar u_n-x|\leq2^{-n}\) by taking the limit;
the source error-to-limit consequence is not a replacement
for its strict pairwise clause.
\end{hbdefinition}

\begin{hbdefinition}{def:p04-fo-native-functions}{Native function names: \(\delta_C\) and \(\delta_{\mathbb R}\)}
Source Definition 6.5 first defines
\(\alpha:\Sigma^*\rightharpoonup\mathrm{Pg}\):
the word \(w=u_0\#v_0\#\cdots\#u_m\#v_m\)
names the linear polygon with vertices
\((\bar u_0,\bar v_0),\ldots,(\bar u_m,\bar v_m)\),
provided every word is a valid rational word and
\(0=\bar u_0<\cdots<\bar u_m=1\).
Then \(\delta_C:\Sigma^\omega\rightharpoonup C[0,1]\)
is defined by
\[
 \begin{split}
 \delta_C(p)=f\quad\text{iff}\quad&
 p=w_0\text{\textcent}w_1\text{\textcent}\cdots,\
 w_n\in\dom(\alpha),\\
 &(\forall n)(\forall j>n)\
 d(\alpha(w_j),\alpha(w_n))<2^{-n},\\
 &(\forall n)\ d(f,\alpha(w_n))\leq2^{-n}.
 \end{split}
\]
Equivalently the last clause can be stated as
\(f=\lim_n\alpha(w_n)\); the strict pairwise clause then
implies the displayed limit-error bounds.

Source Definition 6.1 defines
\(\delta_{\mathbb R}:\Sigma^\omega\rightharpoonup C(\mathbb R)\).
A name has complete records
\[
 p=u_0\#v_0\#x_0\#y_0\text{\textcent}\,
   u_1\#v_1\#x_1\#y_1\text{\textcent}\cdots
\]
of valid rational words and denotes \(f\) exactly when,
for all rational \(a,b,c,d\),
\[
 f([a,b])\subseteq(c,d)
 \quad\Longleftrightarrow\quad
 (\exists i)\ (a,b,c,d)
 =(\bar u_i,\bar v_i,\bar x_i,\bar y_i).
\]
No natural-index padding convention is part of this native
record syntax. Reversed domain intervals are empty;
their vacuous properties are included in the source clause.
\end{hbdefinition}

\begin{hbdefinition}{def:p04-fo-polygons}{Uniform polygon names}
Write \(C[0,1]\) for the real-valued continuous functions on
\([0,1]\), and
\(\lVert f\rVert_\infty=\max_{0\leq x\leq1}\lvert f(x)\rvert\).
A rational polygon has vertices
\((t_0,y_0),\ldots,(t_m,y_m)\), where
\(0=t_0<\cdots<t_m=1\) and all coordinates are rational;
it is linear between successive vertices.
Use an effective finite code recording the number of vertices
and their rational codes from \cref{def:p04-rs-rational-code}.
Valid codes are decidably recognizable. Reserve code zero for
the zero polygon.

The representation \(\gamma_{\mathrm P}\) interprets a Baire
stream of valid polygon codes \(P_n\) as \(f\) exactly when
\[
 \lVert P_n-f\rVert_\infty\leq2^{-n}\quad\text{for every }n.
\]
Only valid names are promised; checking all these inequalities
is not part of the representation.
\end{hbdefinition}

Rational polygons are dense: uniform continuity gives a fine
rational partition on which variation is small, and rational
approximations to its vertex values give a uniformly close
linear interpolant. This proves existence of names, not a
computable modulus for an unnamed continuous function.
The annex makes the construction and uniqueness explicit.

Define the separate host \(\gamma_{\mathrm{pair}}\) on
Baire streams of numerical polygon codes by both the strict
pairwise and limit-error clauses of native \(\delta_C\).
This is a representation bridge, not a new name for the
source-native \(\delta_C:\Sigma^\omega\rightharpoonup C[0,1]\).

\begin{hbdefinition}{def:p04-fo-boxes}{Compact-range information}
A local box-name of \(f\in C[0,1]\) lists exactly all rational
quadruples \((a,b,c,d)\), with
\(0\leq a\leq b\leq1\), \(c<d\), for which
\[
 f([a,b])\subseteq(c,d).
\]
Encode their indices plus one, with zero padding and arbitrary
order/repetition, and call this representation
\(\gamma_{\mathrm B}\).
For \(C(\mathbb R)\), use the same property on every ordered
rational compact interval without restricting \(a,b\);
call it \(\gamma_{\mathbb R}\).
This host index representation is compared with the native
finite-word records by the following theorem.
\end{hbdefinition}

\begin{hbtheorem}{thm:p04-fo-naming}{proved=annex:p04-fo-naming}{Function-name bridges with both directions}
The representations \(\gamma_{\mathrm P}\),
\(\gamma_{\mathrm{pair}}\) and \(\gamma_{\mathrm B}\) are
computably equivalent. Native \(\delta_C\) is computably
equivalent to them by explicit finite-word/code translators;
native \(\delta_{\mathbb R}\) is computably equivalent to
host \(\gamma_{\mathbb R}\).
Native scalar \(\rho_C\) is computably equivalent to the
host scalar \(\delta_{\mathrm C}\).
Restriction from a
\(\gamma_{\mathbb R}\)-named global function to \([0,1]\)
is computable. Conversely, a \(\gamma_{\mathrm P}\)-name
computably gives a global box-name of the constant-endpoint
extension \(x\mapsto f(\min(1,\max(0,x)))\).
Thus the restricted-global representation of source
Theorem 6.6 and the host representation are equivalent.
\end{hbtheorem}

\begin{hbconvention}{def:p04-fo-native-operators}{Native operator input and output representations}
The source evaluation signature is
\((\delta_{\mathbb R},\rho_C,\rho_C)\):
its first operand is a global function name, its second
a scalar argument name, and its output a scalar value name.
Differentiation has signature \((\delta_C,\delta_C)\)
on \(C^1[0,1]\): both operands name functions.
Global integration has signature
\((\delta_{\mathbb R},\rho_C,\rho_C,\rho_C)\),
with two scalar limits and a scalar output;
the unit-interval integral has \((\delta_C,\rho_C)\).
The native zero relation has \((\delta_C,\rho_C)\),
from a named function to a scalar zero.
The host algorithms below are transferred to these native
signatures by the two-way bridges just stated.
\end{hbconvention}

The source's strict pairwise clause needs a precision shift;
mere convergence is not enough. The box-to-polygon proof
searches for an actual finite cover by small-range boxes.
It does not receive a modulus for free.
For a polygon \(P\), its maximum, slope bound, and minimum
of \(\lvert P\rvert\) on a finite union of rational compact
intervals are computable rational data: inspect vertices,
endpoints and any linear zero crossings.

\section{Evaluation, a modulus and integration}
\label{sec:p04-fo-integration}

\begin{hbtheorem}{thm:p04-fo-evaluation}{proved=annex:p04-fo-evaluation}{Evaluation and effective compact control}
Evaluation \( (f,x)\mapsto f(x)\) is computable from
\(\gamma_{\mathrm P}\)-names and scalar Cauchy names,
on the domain \(x\in[0,1]\).
A name also computes a modulus \(\mu\) satisfying
\[
 |x-y|\leq2^{-\mu(n)}
 \ \Longrightarrow\
 |f(x)-f(y)|\leq2^{-n}.
\]
It computes \(\max_{[0,1]}f\) as a scalar real.
The modulus and precision choices may depend on the name;
the evaluated value and maximum do not.
\end{hbtheorem}

The source proves a compact modulus construction in
Theorem 6.4 and Corollary 6.7 and a maximum-value algorithm
in Theorem 6.8. The host proof obtains its modulus directly
from the finite slope bound of an approximating polygon.
This is extra available information in a uniform name;
pointwise values at rationals without a supplied error control
are not being substituted for that name.

\begin{hbtheorem}{thm:p04-fo-integration}{proved=annex:p04-fo-integration}{Integration with scalar and function output}
From a \(\gamma_{\mathrm P}\)-name one can compute
\(\int_0^1f(t)\,dt\) as a scalar and
\[
 Jf(x)=\int_0^x f(t)\,dt
\]
as a \(\gamma_{\mathrm P}\)-named function.
The scalar approximation \(\int_0^1P_n\) is rational and
has error at most \(2^{-n}\).
Integration between any two Cauchy-named limits in
\([0,1]\) is computable, including reversed or equal limits.
For global input \(\gamma_{\mathbb R}\), integration between
arbitrary finite real limits is computable.
\end{hbtheorem}

Source Theorem 6.10 states the global integration result but
omits its proof. The annex supplies exact polygon integrals,
the function-output interpolation error, and the computable
restriction/rescaling for arbitrary real limits.
It neither assumes a global bound on the input function
nor decides equality of the integration limits.
In the native presentation
\(\operatorname{Int}:C(\mathbb R)\times\mathbb R\times
\mathbb R\to\mathbb R\),
\(\operatorname{Int}(f,a,b)=\int_a^bf(x)\,dx\), is
\((\delta_{\mathbb R},\rho_C,\rho_C,\rho_C)\)-computable.
The preceding bridges transfer the supplied construction
to precisely those native operand types.

\begin{hbexample}{def:p04-fo-integral-case}{Value accuracy and area accuracy}
For \(f(x)=x^2\), interpolation on a uniform rational grid
of mesh \(h\) has error at most \(h^2/4\), since on one
cell its excess is \((x-a)(b-x)\).
Its integral approximates \(1/3\) with at most the same
uniform-error bound. The value \(Jf(x)=x^3/3\) is a
function output, not a single numerical area.
A continuous function need not be differentiable for either
integration algorithm to work.
\end{hbexample}

\section{Derivative information is a different promise}
\label{sec:p04-fo-differentiation}

\begin{hbtheorem}{thm:p04-fo-differentiation}{proved=annex:p04-fo-differentiation}{Failure of uniform differentiation}
On the domain \(C^1[0,1]\), the derivative operator
\(f\mapsto f'\), using \(\gamma_{\mathrm P}\) for both
input and output, has no continuous realizer and hence
no computable realizer. Even the scalar operation
\(f\mapsto f'(0)\) has no continuous realizer.
Endpoint derivatives are one-sided, with continuous extension
of the derivative to the closed interval.
\end{hbtheorem}

The native source operator
\(\operatorname{Diff}:C[0,1]\rightharpoonup C[0,1]\),
with domain \(C^1[0,1]\) and value \(f'\), is not
\((\delta_C,\delta_C)\)-continuous (Theorem 6.9).
Both occurrences of \(\delta_C\) name functions here.
The native countersequence in source Theorem 6.9 is
\(f_n(x)=\sin(n\pi x)/n\): function values tend uniformly
to zero, while \(f_n'(0)=\pi\).
The annex also gives a rational-polynomial version
\(h_m(x)=x(1-x)^m\), with uniform size at most
\(1/(m+1)\) and \(h_m'(0)=1\).
A finite low-amplitude function-name prefix can therefore
hide a large derivative at a point.
This is failure of a uniform operator. It is not, by itself,
a proof that one specified computable function has a
noncomputable derivative.

\begin{hbdefinition}{def:p04-fo-derivative-data}{Stronger derivative representations}
A \(C^1\)-pair name consists of names of \(f\) and \(g\)
with the semantic promise \(g=f'\).
Projection of the second name computes differentiation in
this stronger representation.
Alternatively, supply a \(\gamma_{\mathrm P}\)-name of
\(f\in C^1[0,1]\) and a Baire name of a modulus \(\mu'\)
for its derivative:
\[
 |s-t|\leq2^{-\mu'(n)}
 \ \Longrightarrow\ |f'(s)-f'(t)|\leq2^{-n}.
\]
This is additional input information. Neither its validity
nor its computability follows from the ordinary name of \(f\).
\end{hbdefinition}

\begin{hbtheorem}{thm:p04-fo-derivative-promise}{proved=annex:p04-fo-derivative-promise}{Differentiation with an explicit derivative modulus}
Differentiation is computable from the second strengthened
input above. At output position \(n\), choose
\(r\geq\mu'(n+2)\), put \(h=2^{-r}\), and use
\(P_{n+r+3}\) to compute rational forward secants at
the grid points \(0,h,\ldots,1-h\).
Use the last backward secant at \(1\), and interpolate the
secant values linearly on that grid.
The resulting polygon approximates \(f'\) within \(2^{-n}\).
\end{hbtheorem}

\section{Existence, relation solutions and root selection}
\label{sec:p04-fo-zeros}

For a multivalued zero problem a realizer must output a name
of some allowed zero on every promised input name.
It need not select the same zero on two different names
of the same function. A single-valued root selector has
that additional decoded extensionality requirement.
These are different computational questions.

\begin{hbtheorem}{thm:p04-fo-general-zero}{proved=annex:p04-fo-general-zero}{General zero finding is not continuously realizable}
On the domain of functions in \(C[0,1]\) with a nonempty
zero set, the relation asking for any zero has no continuous
realizer. This remains true on the subclass with
\(f(0)<0<f(1)\). Sign change guarantees existence but does
not alone supply a uniform algorithm.
\end{hbtheorem}

The source-native obstruction uses polygons \(G_t\) with
vertices \((0,-1),(1/3,t),(2/3,t),(1,1)\).
At \(t=0\) a whole interval is zero; positive and negative
small perturbations force roots to opposite sides of it.
The annex expresses this as a finite-name-prefix obstruction,
including name-dependent relation outputs
\cite[Theorem 7.2]{Weihrauch1995IntroductionRS}.

\begin{hbdefinition}{def:p04-fo-sign-change-class}{The no-zero-interval promise}
Let \(\mathcal F_{\mathrm{nd}}\) consist of continuous functions
having opposite signs at some two points and not vanishing
identically on any nondegenerate open subinterval of \([0,1]\).
Rational sign witnesses can then be found by strict
semidecision. The no-zero-interval condition is a semantic
promise, not a decidable preliminary test.
\end{hbdefinition}

\begin{hbtheorem}{thm:p04-fo-nonextensional}{proved=annex:p04-fo-nonextensional}{Computable relation without an extensional selector}
For \(\mathcal F_{\mathrm{nd}}\), the relation asking for
any zero is computable, but it has no continuous single-valued
root selector, and hence no computable one.
A positive algorithm finds rational sign brackets and
repeatedly searches for an interior opposite-sign pair of
at most half the previous width, preserving the signs at
the old endpoints. Its fast limit is a zero.
\end{hbtheorem}

The source Theorem 7.3 establishes both clauses.
The annex proves the shrinking-pair search actually succeeds:
it cannot silently stall forever at a zero plateau.
It also reconstructs the source selector counterfamily
with vertices \((0,-1),(1/3,t),(2/3,t-1),(1,1)\).
Its moving roots rule out an extensional continuous choice,
although a name-dependent relation solution is possible.

\begin{hbtheorem}{thm:p04-fo-unique-zero}{proved=annex:p04-fo-unique-zero}{Unique and supplied isolated roots}
On the domain of functions with exactly one zero in
\([0,1]\), that zero is computable from a uniform name.
No sign change, lower slope bound or supplied isolation
modulus is needed. At precision \(n\), search for rational
\(u<v\), \(v-u<2^{1-n}\), and \(k\) such that
\[
 \min_{[0,1]\setminus(u,v)}|P_k|>2^{-k},
\]
and output \((u+v)/2\). An empty complement satisfies the
test vacuously.
An explicitly supplied rational interval containing exactly
one zero gives the analogous local algorithm, even if
other zeros exist outside it. A singleton interval promised
to contain a zero returns its endpoint.
\end{hbtheorem}

Allow \(u<0\) or \(v>1\); then a root at an endpoint is
not accidentally included in the excluded complement.
The source Theorem 7.5 uses a stronger exclusion margin
\(2\cdot2^{-k}\); either positive margin tends to zero and
permits a correct search. The annex verifies the host clause
and explains the source margin separately.
Source Corollary 7.6 gives individual computability of an
isolated zero of a computable function: existence of a rational
isolating interval allows it to be fixed in that one program.
It does not uniformly find such an interval from an arbitrary
function name.

Strictly increasing functions with
\(f(0)\leq0\leq f(1)\) have a unique zero by the intermediate
value theorem, so the unique-root algorithm applies.
Nondecreasing functions can have a zero plateau; uniqueness
must not be inferred from that weaker word.
For global strictly increasing functions onto \(\mathbb R\),
the same bracketing principle computes the inverse:

\begin{hbtheorem}{thm:p04-fo-inverse}{proved=annex:p04-fo-inverse}{A promised global monotone inverse}
On the domain of \(\gamma_{\mathbb R}\)-named strictly
increasing functions with range \(\mathbb R\), both
\((f,y)\mapsto f^{-1}(y)\) and the inverse function as a
\(\gamma_{\mathbb R}\)-output are computable.
Strictness and surjectivity are promises, not established
by sampling finitely many values.
\end{hbtheorem}

\begin{hbtheorem}{thm:p04-fo-approximate-zero}{proved=annex:p04-fo-approximate-zero}{Residual error is not root-location error}
Given a name of \(f\) and \(n\), one can find rational
\(s\in[0,1]\) with \(|f(s)|<2^{-n}\) whenever such a
point exists. Search for a certificate
\[
 |P_k(s)|+2^{-k}<2^{-n}.
\]
In particular, the search succeeds when \(f\) has a zero.
Without additional separation information this supplies
no bound on the distance from \(s\) to a zero.
\end{hbtheorem}

The printed search in source Theorem 7.7 lacks the precision
control needed to transfer a small polygon value to the
true function. The displayed certificate supplies it.
This is an explicit corrected host algorithm, not a claim
that the source printed the same test.

\begin{hbexample}{def:p04-fo-zero-cases}{Four different zero questions}
The function \(f(x)=x\) has its unique root at an endpoint.
The function \((x-1/2)^2\) has a unique interior root but
does not change sign; the unique-root exclusion algorithm
still applies. A three-root polynomial
\((x-1/4)(x-1/2)(x-3/4)\) permits different legitimate
relation answers; an isolating interval is extra data.
The constant function \(f(x)=\epsilon>0\), with
\(\epsilon<2^{-n}\), has approximate residuals everywhere
but no exact root. Small residual is not an existence proof.
For the zero function every point is a root: it lies in the
general relation domain but neither the unique-root nor
no-zero-interval domain.
\end{hbexample}

\begin{hbtheorem}{thm:p04-fo-no-computable-zero}{proved=annex:p04-fo-no-computable-zero}{Existence need not include a computable zero}
There is a computable continuous function on \([0,1]\)
with a nonempty zero set and no computable zero.
It follows that there is a computable continuous function
with a computable maximum value but no computable maximum
point.
\end{hbtheorem}

The proof constructs an effectively enumerated open cover
of all computable reals with a small total interval-length
budget, then uses the source's weighted nonnegative tent
functions to make its complement the zero set
\cite[Theorem 7.1 and following cited boundary]{Weihrauch1995IntroductionRS}.
The source cites an external Specker construction rather
than giving that construction in full; the annex supplies
a direct small-cover and tent-function proof.

\begin{proofboundary}{def:p04-fo-source-boundary}{Declared extent}
All displayed positive algorithms and obstructions have full
local annex proofs under the named representations and promises.
The 1995 universal characterization of evaluation, its broader
set/image/composition results and its source-admissibility
assertion retain their original omitted/cited boundaries unless
explicitly reconstructed here.
The constructed algorithms establish only their stated
representations and promises; they do not identify
topological continuity with computability.
\end{proofboundary}
